% GENERATED by scripts/build.py convert abstract -- do not hand-edit; edit
% chapters/md/08_abstract.md and regenerate.
% md sha256 f6878af622d52c0fbc772c8416d040e2d0c06a256b68dff7a3dfcec44b7013d8
% canonical manuscript/v13_2026-10-17/08_abstract.md sha256 f6878af622d52c0fbc772c8416d040e2d0c06a256b68dff7a3dfcec44b7013d8
% md identical to canonical: True
% filter sha256 396c91097c881993bcd2e433e51e3f72c32ab14e820fa613277c56b07f4b3ca9

\section*{Abstract}\label{sec:abstract}\addcontentsline{toc}{section}{Abstract}

This thesis asks whether water and sanitation (WASH) conditions are associated with the anthropometric outcomes of children under five, measured as height-for-age (HAZ), weight-for-age (WAZ) and weight-for-height (WHZ) z-scores, in four Sub-Saharan African countries: Ethiopia, Ghana, Kenya and Nigeria. The baseline analysis distinguishes between a child's own household facilities and the WASH conditions of the surrounding community. The four countries were selected for data availability and recency, not to represent the region as a whole. Both analyses are cross-sectional and observational, and every estimate is reported as an association, not as a causal effect.

The baseline sample pools the most recent Demographic and Health Survey (DHS) round for each country (2022 to 2025). Household WASH is the improved or unimproved facility classification of the child's own household. Community WASH is a leave-one-out rate of improved-facility coverage among the other sampled households in the same DHS cluster, so a child's own household never enters its community measure. Each exposure is estimated under the fixed-effects structure its construction requires: cluster fixed effects for household WASH, which compare households within the same cluster, and region fixed effects for community WASH, which varies almost only between clusters. No household cluster-fixed-effects coefficient reaches significance at the 5\% level for any of the three outcomes. Household water and HAZ comes closest (\emph{p}~=~.062, crossing only a 10\% threshold); for sanitation and HAZ, \emph{p}~=~.548. These results are inconclusive rather than evidence that no association exists. Community coverage shows a large, significant association before covariate adjustment for most outcome and exposure pairs; the exception is water and WHZ, which is small and non-significant from the start. Once socioeconomic controls and region fixed effects are added, none of the community associations remains statistically significant, and the fully adjusted HAZ estimates for water (\emph{p}~=~.573) and sanitation (\emph{p}~=~.894) are both inconclusive. This attenuation is consistent with confounding by local socioeconomic conditions, but it does not establish confounding as the explanation.

A separate, preliminary extension uses an independently constructed sample: earlier DHS rounds for Ghana (2014), Kenya (2014) and Nigeria (2018), harmonised by IPUMS-DHS and linked to modelled annual gridded estimates of local improved water and sanitation coverage from the Institute for Health Metrics and Evaluation (IHME). It responds to one limitation of the baseline design, which records WASH once, at the survey date, rather than during each child's own early life. Within this sample, it asks whether modelled coverage during a child's own early-life window, here the first 12 postnatal months, is associated with later anthropometric outcomes, using cluster fixed effects and country-specific birth-year terms. Ethiopia's 2016 round enters only as a provisional sensitivity analysis, because its interview-date conversion could not be validated. The exposure is an area average of modelled coverage around displaced cluster coordinates; it is neither observed household access nor population-weighted coverage. Because the historical data do not separate a child's household from its community, the extension does not repeat the household-versus-community comparison.

The selected primary finding is the pooled established-date association between local water coverage and HAZ: β = 1.03 per 10 percentage points of modelled coverage (95\%~CI {[}0.51,~1.54{]}, \emph{p}~\textless~.001). It is one of six nominally significant cells (\emph{p}~\textless~.05) among the 24 primary cells, and no adjustment for multiple testing was applied. The estimate depends on the equal-country weighting convention: unweighted, on the same 23,018 observations, it is 0.41 (95\%~CI {[}0.03,~0.80{]}, \emph{p}~=~.035), and equal country weight totals do not give the three countries equal influence on the pooled slope. All 24 primary cells were reproduced with a second estimation library on a separately rebuilt sample; the other 267 of the 291 historical models passed internal checks and, apart from two unweighted counterparts, were not re-estimated with a second library. This reproduction confirms the computation, not causality or robustness to multiple testing. The models are weighted with the IPUMS-DHS PERWEIGHT variable. No official documentation of an anthropometry-specific weight for the measured subsample could be located, so the treatment of measurement selection rests on a stated working inference, and the estimates are weighted conditional associations among eligible measured children rather than population-representative estimates. The exposure source also draws on the Ghanaian and Kenyan DHS surveys themselves, and current residence stands in for early-life residence.

The baseline and historical samples differ at the same time in their children and survey rounds, exposure construction, eligibility rules and control sets, so the contrast between the largely null baseline results and the historical finding cannot be attributed to exposure timing. The historical result is a preliminary, qualified finding and does not overturn the more cautious baseline conclusions. The thesis contributes a side-by-side comparison of household and leave-one-out community WASH within one four-country sample and a preliminary, qualified extension that assigns modelled local coverage to each child's early-life window, each read with its own limitations. None of the results is evidence for or against the effectiveness of WASH investment as a policy instrument.
